% =============================================================================
% TECHNICAL WHITEPAPER TEMPLATE
% 40-60 páginas, estilo industrial
% =============================================================================

% Usar el preámbulo base
% =============================================================================
% PREÁMBULO LATEX EXACTO - COPIAR LITERALMENTE
% =============================================================================
\documentclass[12pt,a4paper]{report}

% Fuentes CRÍTICAS para compatibilidad Word
\usepackage[utf8]{inputenc}
\usepackage[T1]{fontenc}
\usepackage{helvet}           % Helvetica para texto principal
\usepackage{courier}           % Courier para código
\usepackage{graphicx}
\usepackage{fancyhdr}
\usepackage{hyperref}
\usepackage{listings}
\usepackage{xcolor}
\usepackage{booktabs}
\usepackage{geometry}
\usepackage{microtype}
\usepackage{tocbibind}
\usepackage{amsmath}
\usepackage{amssymb}
\usepackage{appendix}
\usepackage{caption}
\usepackage{subcaption}
\usepackage{enumitem}
\usepackage{titlesec}

% Márgenes
\geometry{
    left=2.5cm,
    right=2.5cm,
    top=2.5cm,
    bottom=2.5cm,
    headheight=15pt
}

% Listings para código
\lstset{
    basicstyle=\ttfamily\fontfamily{pcr}\selectfont,
    backgroundcolor=\color{gray!10},
    frame=single,
    frameround=tttt,
    numbers=left,
    numberstyle=\tiny\color{gray},
    keywordstyle=\color{blue}\bfseries,
    commentstyle=\color{green!60!black},
    stringstyle=\color{red},
    breaklines=true,
    tabsize=2,
    captionpos=b,
    xleftmargin=2em,
}

% Headers y footers
\pagestyle{fancy}
\fancyhf{}
\fancyhead[L]{\color{blue!70!black}\textbf{\sffamily\leftmark}}
\fancyhead[R]{\color{blue!70!black}\textbf{\sffamily\thepage}}
\fancyfoot[C]{\color{gray}\small\sffamily Project Name}

% Títulos con sans-serif
\titleformat{\chapter}[display]
    {\normalfont\huge\bfseries\sffamily\color{blue!70!black}}
    {\chaptertitlename\ \thechapter}{20pt}{\Huge}

\titleformat{\section}
    {\normalfont\Large\bfseries\sffamily\color{blue!70!black}}
    {\thesection}{1em}{}

\titleformat{\subsection}
    {\normalfont\large\bfseries\sffamily\color{orange!80!black}}
    {\thesubsection}{1em}{}

% Hyperref
\hypersetup{
    colorlinks=true,
    linkcolor=blue!70!black,
    urlcolor=blue!70!black,
    pdftitle={Project Name},
    pdfauthor={wisrovi},
}

% =============================================================================
% REGLAS DE ORO
% =============================================================================
% 1. USAR SOLO helvet (\sffamily) + courier (\ttfamily)
% 2. EVITAR lmodern, fuentes matemáticas complejas
% 3. CAPÍTULOS: \cleardoublepage al inicio
% 4. TÍTULOS: \sffamily en todos
% 5. CÓDIGO: \fontfamily{pcr}\selectfont
% 6. IMÁGENES: Solo PNG/JPG, rutas a sources/


\begin{document}

% =============================================================================
% PÁGINA DE TÍTULO
% =============================================================================
\begin{titlepage}
    \centering
    \vspace*{2cm}
    
    \ifdefined\sourcespath
    \includegraphics[width=0.25\textwidth]{\sourcespath/project-logo.png}
    \fi
    
    \vspace{1cm}
    
    {\Huge\textbf{\sffamily Project Title}}\\[0.5cm]
    {\Large\textbf{\sffamily Subtitle}}\\[1cm]
    {\large\textbf{\sffamily \today}}
    
    \vfill
    
    \colorbox{blue!70!black}{\color{white}\textbf{\sffamily Technology Stack}}
    
    \vspace{2cm}
    
    {\large\textbf{\sffamily Technical Whitepaper}}\\[0.5cm]
    {\large\textbf{\sffamily Document Version 1.0}}
\end{titlepage}

% =============================================================================
% PÁGINA DE AUTOR
% =============================================================================
\begin{titlepage}
    \centering
    \vspace*{3cm}
    
    {\Large\textbf{\sffamily About the Author}}\\[1cm]
    {\Large\textbf{\sffamily wisrovi}}\\[0.5cm]
    \textit{\sffamily AI Solutions Architect}\\[1cm]
    
    \vspace{1cm}
    
    \textbf{\sffamily Contact Information:}\\[0.5cm]
    \href{https://www.linkedin.com/in/wisrovi-rodriguez}{\sffamily LinkedIn}\\[0.3cm]
    \href{https://github.com/wisrovi}{\sffamily GitHub}\\[0.3cm]
    \href{https://hub.docker.com/u/wisrovi}{\sffamily DockerHub}\\[0.3cm]
    \href{https://pypi.org/search/?q=wisrovi}{\sffamily PyPI}\\[0.3cm]
    \hrefmailto{wisrovi.rodriguez@gmail.com}{\sffamily wisrovi.rodriguez@gmail.com}
\end{titlepage}

% =============================================================================
% TABLA DE CONTENIDOS
% =============================================================================
\tableofcontents
\cleardoublepage
\listoffigures
\cleardoublepage
\listoftables
\cleardoublepage

% =============================================================================
% CAPÍTULOS
% =============================================================================

\chapter{Executive Summary}
\cleardoublepage

\section{Overview}
% Breve descripción del proyecto

\section{Key Findings}
% Principales hallazgos técnicos

\section{Recommendations}
% Recomendaciones de alto nivel

% =============================================================================

\chapter{Introduction}
\cleardoublepage

\section{Background}
% Contexto y motivación

\section{Objectives}
% Objetivos del documento

\section{Scope}
% Alcance del documento

% =============================================================================

\chapter{Architecture}
\cleardoublepage

\section{System Overview}
% Vista general del sistema

\section{System Architecture}
\begin{figure}[htbp]
    \centering
    \includegraphics[width=0.95\textwidth]{sources/system-architecture.png}
    \caption{System Architecture - High-level component overview}
    \label{fig:arch-overview}
\end{figure}
\textit{Figure \ref{fig:arch-overview} illustrates the primary system components and their interactions.}

\section{Data Flow}
\begin{figure}[htbp]
    \centering
    \includegraphics[width=0.95\textwidth]{sources/data-flow.png}
    \caption{Data Flow - End-to-end data transformation}
    \label{fig:data-flow}
\end{figure}
\textit{Figure \ref{fig:data-flow} demonstrates the complete data pipeline from input to output.}

\section{Component Interaction}
\begin{figure}[htbp]
    \centering
    \includegraphics[width=0.95\textwidth]{sources/component-interaction.png}
    \caption{Component Interaction - Module communication patterns}
    \label{fig:component-interaction}
\end{figure}
\textit{Figure \ref{fig:component-interaction} details the inter-module communication protocols.}

% =============================================================================

\chapter{Technical Analysis}
\cleardoublepage

\section{Code Quality Metrics}
% Métricas de calidad de código

\section{Security Assessment}
% Evaluación de seguridad

\section{Performance Considerations}
% Consideraciones de rendimiento

% =============================================================================

\chapter{Risk Assessment}
\cleardoublepage

\section{Technical Risk Matrix}

\begin{table}[htbp]
\centering
\caption{Technical Risk Matrix}
\label{tbl:risk-matrix}
\begin{tabular}{|l|c|c|c|p{4cm}|}
\hline
\textbf{Component} & \textbf{Risk} & \textbf{Prob} & \textbf{Impact} & \textbf{Mitigation} \\
\hline
Database & High & 3 & 5 & Connection pooling \\
\hline
API Rate Limits & Medium & 4 & 3 & Exponential backoff \\
\hline
Data Loss & Medium & 2 & 4 & Regular backups \\
\hline
\end{tabular}
\end{table}

\section{Mitigation Strategies}
% Estrategias de mitigación

% =============================================================================

\chapter{Conclusion}
\cleardoublepage

\section{Summary}
% Resumen de hallazgos

\section{Future Work}
% Trabajo futuro recomendado

% =============================================================================

\appendix
\chapter{Glossary}
% Glosario de términos

\chapter{Bibliography}
% Referencias en formato APA

\begin{thebibliography}{99}
%\bibitem{key} Author. (Year). Title. Publisher.
\end{thebibliography}

\end{document}
